\scanpage{9}

\mspage{5}

d'un élément du groupe de Néron Severi de $J^m$) est \uline{positive} i.e. l'élément en question de Néron--Severi est ample, i.e. une polarisation.

Pratiquement, cela s'explicite ainsi: Soit $T$ une variété de paramètres \inserted{connexe} non singulière munie d'un point marqué $a$, $z$ une classe de cycles de codimension $m$ sur $T\mathbin{x}X$ (à équivalence linéaire près, mettons), telle que $z(a)=0$ dans $X$, soi\unsure{e|t}nt $p$ et $q$ les deux projections de $T\mathbin{x}T\mathbin{x}X$ sur $T\mathbin{x}X$, $r$ sa projection sur $T\mathbin{x}T$, considérons
\[
D=r_*(p^*(z)\mathbin{.}q^*(z))
\]
qui est une classe de diviseurs sur $T\mathbin{x}T$, que nous considérons comme une classe de correspondance divisorielle sur $T\mathbin{x}T$. Si $A=\operatorname{Alb}^0(T)$ (NB si tu veux, tu peux supposer $T=A$ et $a$ l'origine), elle provient donc d'une classe de correspondance sur $A\mathbin{x}A$, évidemment symétrique. Soit $N$ le ``noyau'' de cette classe (i.e. le noyau de $A\longrightarrow(\text{duale de }A)$ qu'elle définit), on obtient alors une classe de correspondance symétrique sur $J\mathbin{x}J$, où $J=A/N$. A prouver que cette dernière est \uline{positive} ! Je me demande si les spécialistes ``abéliens'' pourraient avoir une idée sur une telle question, peut-être Matsusaka ? Ou toi-même ? Notes d'ailleurs que cette question te dévoile pratiquement la méthode de construction des $J^i$ généraux; Si tu veux, tu peux te borner \inserted{aussi} au cas où tu disposes d'une sous-variété de codimension $m-1$ $Y$ de $X$, non singulière si tu y tiens, \deleted{\unsure{tu|xx}} et où $T=\operatorname{Pic}^0(Y)$, considéré comme paramétrant les diviseurs alg équiv à zéro de $Y$, mais considérés comme classes de cycles \deleted{dans} de codimension $m$ de $X$.

Merci pour la copie de la lettre à Ogg !

Bien à toi

\mspage{8}
